% A Research Letter draft for Advances in Radiation Oncology, modeled on
% Capaldi et al. 2026 (https://doi.org/10.1016/j.adro.2025.101970).
%
% Build with:  pdflatex main && bibtex main && pdflatex main && pdflatex main
%
% Author order is a placeholder.  The clinical leadership question is open;
% see TODO.md.

\documentclass[3p,times]{elsarticle}

\usepackage{graphicx}
\usepackage{amsmath}
\usepackage{amssymb}
\usepackage{booktabs}
\usepackage{url}
\usepackage{hyperref}
\hypersetup{colorlinks=true,linkcolor=blue,citecolor=blue,urlcolor=blue}

\journal{Advances in Radiation Oncology}

\begin{document}

\begin{frontmatter}

\title{Tide: A Browser-Based At-Home Coaching System for Deep-Inspiration
Breath-Hold Reproducibility in Left-Sided Breast Radiation Therapy}

\author[mdacc]{Ramez Kouzy}
\ead{rkouzy@mdanderson.org}

\author[mdacc]{Melissa P.~Mitchell}
\author[breath]{Belisa Vranich}
\author[mdacc]{(other co-authors TBD)}

\address[mdacc]{Department of Radiation Oncology, The University of Texas
MD Anderson Cancer Center, Houston, Texas, USA}
\address[breath]{The Breathing Class, New York, NY, USA}

\cortext[cor1]{Corresponding author.}

\begin{abstract}
\textbf{Purpose.} Deep-inspiration breath-hold (DIBH) reduces cardiac
dose during left-sided breast radiation therapy, but its dosimetric
benefit depends on the patient's ability to reproduce a consistent
chest-hold position across fractions. Surface-guided radiation therapy
(SGRT) handles this in the treatment vault; no comparable feedback is
available at home, where between-sim and between-fraction practice
could plausibly improve preparedness and on-couch reproducibility. We
describe \emph{Tide}, a browser-based smartphone application that uses
the phone's fused-attitude orientation sensor to coach DIBH
reproducibility at home, with adaptive per-session calibration and
position-match audio feedback.

\textbf{Methods and Materials.} Tide was implemented as a static
single-page application (Next.js 16, React 19) deployed on Vercel,
delivering identical behaviour on iOS Safari and Android Chrome via
\texttt{DeviceOrientationEvent}. The patient places the phone in
portrait with the charging-port edge resting on the sternum, allowing
the phone to act as a lever during chest expansion. A guided
four-step session (Calibration $\rightarrow$ Learn $\rightarrow$
Practice $\rightarrow$ Report) measures (i) the patient's
breathing-pitch standard deviation over 12\,s of normal breathing,
(ii) a target plateau pitch from the average of three ``comfortable''
calibration holds, and (iii) the reproducibility and on-target time
of five subsequent practice holds. In-hold stability is detected
adaptively via a 2-second rolling SD threshold set to
$0.7\times$ baseline breathing SD, with a hysteretic state machine
(1.0\,s lock debounce, 1.5\,s drift debounce). Coaching audio is
delivered as pre-rendered MP3 clips (ElevenLabs, ``Rachel'' voice)
through a single \texttt{<audio>} element; the Web Speech API is the
silent fallback if a clip is missing. Haptic confirmation uses the
Web Vibration API where supported (Android only).

\textbf{Results.} Two illustrative recordings (single subject, Apple
iPhone, iOS Safari) demonstrate the design. With the phone laid flat
and centered on the abdomen (belly-only breathing), normal-breathing
pitch SD was 0.81$^{\circ}$ and in-hold pitch SD 0.07$^{\circ}$ (11$\times$
ratio); peak inhale pitch deviation was small (5$^{\circ}$). With the
phone bridging from sternum to upper belly (chest-led DIBH-style
breathing), normal-breathing SD was 0.99$^{\circ}$ and in-hold SD
0.13$^{\circ}$ (6$\times$ ratio); peak inhale pitch deviation rose to
14$^{\circ}$. The sternum-bridging placement therefore yielded
substantially larger absolute signal at modest cost in in-hold
quietness. The adaptive threshold cleanly separated breathing from
holds in both modes.

\textbf{Conclusions.} A consumer smartphone, delivered through a
browser without any add-on hardware, can capture sufficient signal
for at-home DIBH reproducibility coaching. Tide complements rather
than replaces in-clinic SGRT systems by addressing the between-fraction
practice gap. Forthcoming work will validate Tide in a single-arm
feasibility cohort against simulator-day plateau reproducibility.
\end{abstract}

\begin{keyword}
deep-inspiration breath-hold \sep DIBH \sep mobile health \sep
breast radiation therapy \sep cardiac sparing \sep
patient-reported outcomes \sep smartphone sensors
\end{keyword}

\end{frontmatter}

\section{Introduction}

DIBH is now a standard-of-care technique for left-sided breast
radiation therapy because it reduces mean heart dose by displacing
the heart away from the chest wall during chest
expansion~\cite{darby2013ischemic,piroth2019heart}. Its dosimetric
benefit, however, depends on the patient reproducing the same chest
position across each treatment fraction; small departures from the
planning position translate directly to dose-distribution
errors~\cite{tang2015dosimetric}. SGRT systems handle this within
the treatment vault, providing real-time positional feedback and
beam-hold gating with sub-millimetre accuracy. Recent work has
demonstrated that a smartphone's onboard sensors are themselves
sufficient to act as an SGRT-equivalent in-room monitor: Capaldi et
al.\ first validated an iOS accelerometer/gyroscope-based monitor
(``IRF'') against the commercial Real-time Position Management
(RPM) system in 2020~\cite{capaldi2020}, and extended this in 2026
with a LiDAR-based six-degree-of-freedom system (``iSGRT'')
validated against in-vault clinical SGRT~\cite{capaldi2026}.

A separate body of work has examined the role of patient
\emph{preparation} for DIBH in the days before simulation and
treatment. Educational mobile interventions reduce
sim-day anxiety and self-reported preparedness without
quantitatively coaching the breath-hold itself~\cite{mitchell2025prep,
dibhapp}. Spirometry-based at-home preparation devices have been
piloted, but their adoption is limited by the cost and complexity
of dedicated hardware~\cite{delombaerde2021spirometer}.

The gap between these two literatures is the
\emph{between-fraction at-home practice} of the breath-hold itself
with quantitative feedback that does not require the patient to
visit the treatment vault. Mitchell et al.\ articulated this gap
directly in early 2026, asking whether a phone laid on the abdomen
during DIBH could (a) measure the depth of inhale, and (b) coach
the patient back to the same plateau position in subsequent
attempts. Tide is our response. It uses the same fused-attitude
sensor surface that Capaldi's IRF validated, but is positioned for
home use, framed as a wellness/practice tool rather than a clinical
motion monitor, and explicitly optimised for
\emph{reproducibility} of the chest-hold position rather than
absolute positional accuracy.

\section{Methods and Materials}

\subsection{Application architecture}

Tide is a single-page client-side application built with Next.js
16.2 (Turbopack) and React 19, deployed as static assets to
Vercel's CDN. There is no backend; all sensor processing, state
machine, and audio playback run in the browser. Sensor input is
read through \texttt{DeviceOrientationEvent} (fused pitch / roll /
yaw) and \texttt{DeviceMotionEvent} (raw acceleration, acceleration
including gravity, and rotation rate). A \texttt{WakeLock} request
is issued on entry to the active phases to prevent the screen from
sleeping during recording.

\subsection{Phone placement and signal model}

The patient lies supine and places the phone in portrait
orientation with the charging-port edge anchored on the sternum,
the rest of the phone extending caudally across the upper abdomen.
This placement makes the phone act as a lever across the
chest-belly seam: chest expansion lifts the cranial end of the
phone while the caudal end (resting on more compliant abdominal
tissue) lags, producing rotation around the medio-lateral axis.
The fused-attitude beta angle reads this rotation directly. A
flat, centered phone produces only translation (both ends rise
together) and no usable rotational signal, as confirmed in our
preliminary recordings (Section~\ref{sec:results}).

\subsection{Four-phase session}

Each session walks the patient through:

\paragraph{Phase 1 --- Calibration (12\,s).}
Twelve seconds of normal breathing, with the first two seconds
discarded as a settling window. The remaining ten seconds yield
the patient's breathing-pitch standard deviation
(\texttt{breathingSD}), mean baseline pitch, and an estimated
breathing rate from zero-crossings of the centered signal. Sanity
gates reject calibrations with implausible amplitude
($<\!0.4^{\circ}$ or $>\!30^{\circ}$) or breath rate (outside
4--30 breaths per minute), prompting the patient to retry.

\paragraph{Phase 2 --- Learn (3 comfortable holds).}
The patient performs three ``comfortable'' chest holds, each at a
depth they can reproduce. A hold is delineated by a Start tap; the
algorithm watches stability passively. Each hold's plateau is
defined as the median pitch over the last 3 seconds of its longest
stable run. After three holds, the session target and tolerance
are derived:
\begin{align}
\text{target} &= \text{mean}(\text{plateau}_{1..3}) \\
\text{tolerance} &= \max(0.5^{\circ},\; 2 \times \text{SD}(\text{plateau}_{1..3}))
\end{align}
The 0.5$^{\circ}$ floor prevents a coincidentally tight calibration
from making subsequent practice impossible; the SD-based scaling
gives more variable patients a more forgiving band.

\paragraph{Phase 3 --- Practice (5 holds with position-match coaching).}
The patient is coached to reproduce the target across five
practice holds. During each hold the algorithm tracks two signals:
\emph{stability} (in-hold rolling SD below an adaptive threshold)
and \emph{position} (current pitch within $\pm$tolerance of
target). Audio cues fire only on stable transitions:
``a little deeper'' (below target), ``ease back a touch'' (above
target), or ``right there, hold steady'' (on target). The user
interface includes a dashed target ring rendered around an animated
breath orb at the size the orb should reach when the patient is on
target; the patient inhales until the live orb fills the ring.
Time-on-target accumulates toward the duration goal (default 20
seconds, configurable 15--35).

\paragraph{Phase 4 --- Report.}
Per-session metrics include: per-hold longest stable run, longest
on-target run, plateau pitch, plateau SD, drift events,
time-to-lock, and total duration. The headline clinical metric is
the standard deviation of plateau pitch across the five practice
holds --- the patient's intra-session reproducibility.

\subsection{Adaptive stability detection}

In-hold stability is determined every 100\,ms by comparing the
2-second rolling SD of pitch to an adaptive threshold:
\begin{equation}
\theta_{\text{stable}} = \mathrm{clamp}\!\left(0.7 \times \mathrm{SD}_{\text{breathing}},\; 0.08^{\circ},\; 1.2^{\circ}\right)
\end{equation}
The 0.7 multiplier was chosen to leave clean holds comfortably
below threshold across the two breathing modes observed
(Section~\ref{sec:results}). A hysteretic state machine prevents
flapping: a transition into ``stable'' must persist for 1.0\,s
before a \emph{lock} event fires; a transition out must persist
for 1.5\,s before \emph{drift} fires. The first lock during a hold
emits the \texttt{locked\_in} cue with a haptic confirmation buzz;
subsequent locks after a drift emit \texttt{regained}.

\subsection{Auto-release}

Two paths automatically end a hold. First, once the patient has
locked at least once, a sustained SD spike above
$3\times\theta_{\text{stable}}$ for 1.5\,s ends the hold; this
detects exhalation as a movement event rather than as an absolute
pitch change, which proved unreliable for low-amplitude patients
(see Section~\ref{sec:results}). Second, a 5-second grace timer
fires after \texttt{target\_reached} so that successful holds end
gracefully rather than running indefinitely.

\subsection{Audio coaching}

The coaching script comprises 26 fixed phrases covering
calibration, Learn-phase progression, stability events
(locked / drifting / regained), position cues
(deeper / ease / right-there), and session bookends. To eliminate
runtime synthesis cost and latency, each phrase is rendered once
to MP3 via the ElevenLabs API (model
\texttt{eleven\_multilingual\_v2}, voice ``Rachel'') and served
from \texttt{public/audio/}. A \texttt{playClip(key)} helper pauses
any in-flight clip and starts the requested one; if the file is
missing or autoplay is blocked, the helper falls back to the
browser's \texttt{SpeechSynthesis} API for the same text. Haptic
confirmations use \texttt{navigator.vibrate}, supported on Android
Chrome but not iOS Safari.

\section{Results}\label{sec:results}

We recorded two illustrative sessions on an Apple iPhone running
iOS 18 / Safari, with the multi-channel lab tool
(\texttt{/lab}) capturing all 13 sensor channels at $\approx$58\,Hz.
The two sessions were performed by a single subject and differed
only in breathing strategy and phone placement.

\paragraph{Belly-mode (phone centered, abdominal breathing).}
The phone was laid flat in landscape with its long axis across the
abdomen at the level of the umbilicus; the subject deliberately
held their chest tight and inhaled diaphragmatically. Across two
holds:
\begin{itemize}
  \item Normal-breathing pitch SD: 0.81$^{\circ}$;\, in-hold SD: 0.07$^{\circ}$ (ratio 11$\times$).
  \item Pitch range during inhale: $\sim\!5^{\circ}$.
  \item Hold mean pitch was indistinguishable from baseline mean ($-1.5^{\circ}$ vs.\ $-1.5^{\circ}$).
\end{itemize}

The phone here acted essentially as a passive translation sensor
rather than a lever: bilateral abdominal rise lifted both ends
together, and the rotational signal during a deep inhale was
trivial. Stability detection still worked, since the patient was
quiet enough during the hold that the ratio was large, but
absolute pitch carried no information about depth and the breath
orb visualisation was static.

\paragraph{Chest-mode (phone bridging sternum to upper belly,
DIBH-style breathing).} The phone was placed in portrait with the
charging-port edge anchored on the sternum, the rest extending
caudally; the subject performed full chest-led DIBH-style holds.
Across one guided 20-second hold and one extended free-form hold:
\begin{itemize}
  \item Normal-breathing pitch SD: 0.99$^{\circ}$;\, in-hold SD: 0.13$^{\circ}$ (ratio 6$\times$).
  \item Pitch range during inhale: $\sim\!14^{\circ}$.
  \item Adaptive threshold $\theta_{\text{stable}} = 0.69^{\circ}$ comfortably separated the in-hold SD (0.13$^{\circ}$) from breathing SD (0.99$^{\circ}$).
\end{itemize}

The sternum-bridging placement transformed the signal: the chest
expansion produced a measurable several-degree pitch swing that
drove the breath orb visualisation, and the in-hold SD remained
small enough that the algorithm cleanly locked. The 6$\times$
hold-to-breathing SD ratio is smaller than the 11$\times$ of
belly-mode (the chest never fully comes to rest the way a tight
chest does), but is far more clinically appropriate: chest-led
DIBH is what produces the cardiac displacement on which the entire
technique depends.

\paragraph{Algorithm validation.} A simulation of the production
algorithm against the chest-mode recording correctly identified the
target-setting hold's stable plateau (15.8\,s of stable run within
a 19.5\,s hold), correctly classified a partial hold (1.4\,s
stable within 6.8\,s of effort) as a non-attempt, and correctly
flagged a single in-hold drift event during a long held attempt.
Subjective testing of the position-match coaching is in progress;
a planned single-arm feasibility study will quantify
inter-session plateau drift in 30--50 left-breast DIBH patients
between sim and the first treatment fraction.

\section{Discussion}

\paragraph{Differentiation from prior work.} Tide deliberately occupies
a different point in the design space from existing
smartphone-based DIBH tools. Capaldi 2020 (IRF) and Capaldi 2026
(iSGRT) target in-vault SGRT replacement: the patient is on the
treatment table, the phone is mounted on a precision arm, and
sub-millimetre accuracy is required~\cite{capaldi2020,capaldi2026}.
The DIBH-app (Sweden) and the Mitchell preparatory app are
educational interventions designed to reduce pre-simulation
anxiety and improve patient understanding, but do not
quantitatively coach the breath-hold
itself~\cite{dibhapp,mitchell2025prep}. Tide is patient-facing and
quantitative, but designed for at-home practice between sim and
treatment, where reproducibility of the patient's
\emph{personal target} matters far more than absolute positional
accuracy. The choice of a browser delivery model (no App Store
review, single codebase across iOS and Android, doctor-texts-a-URL
distribution) is a direct reflection of this design constraint.

\paragraph{Adaptive thresholds.} Our preliminary recordings make
clear that no fixed stability threshold can serve a population:
breathing-pitch SD differs by an order of magnitude across
breathing strategies, and would differ further across body
habitus, phone placement, and patient compliance. The adaptive
threshold ($0.7 \times \text{breathingSD}$) was specifically
chosen to leave both observed modes a comfortable margin below
threshold. Wider validation may motivate further refinement, but
the principle of per-session adaptation is, we believe, essential.

\paragraph{Stability versus position.} A common failure mode in
early prototypes was rewarding ``hold still anywhere''. The Learn
phase, by averaging three comfortable holds, defines a
patient-specific target plateau, and Practice rewards only the
intersection of stability and position-match. This corresponds
directly to the clinical reality that DIBH cardiac sparing
requires a particular chest position, not just any held position.

\paragraph{Limitations.} Three caveats merit emphasis. First, the
results above are from a single subject across two breathing
modes; population-level validation is essential and is the subject
of forthcoming work. Second, Tide does not measure absolute lung
volume or absolute chest position; it measures relative pitch and
is calibrated per-session, so cross-session drift is currently
unobservable. (A localStorage-based cross-session anchor is
deferred future work; see TODO.md.) Third, the application is
intentionally framed as a wellness / practice tool, not a Software
as a Medical Device; any cross-over to clinical decision support
would require a 510(k) regulatory pathway, likely with iSGRT or
RPM as the predicate device. Fourth, the patent landscape includes
Capaldi et al.'s pending Surface Audio-Visual Biofeedback (SAVB)
patent (assigned to UCSF / Stanford), whose claim language we have
not yet reviewed in detail; Tide's home-use framing and absence of
LiDAR may render this moot, but a formal review is required before
any commercial path.

\paragraph{Forthcoming work.}
Planned next steps include (i) a single-arm IRB-approved feasibility
study at MD Anderson Cancer Center comparing intra-session
plateau-pitch SD to a same-day simulator measurement; (ii)
re-recording of the coaching audio with Belisa Vranich's voice, in
keeping with her existing clinical-psychology breathing
instruction methodology; (iii) outreach to the Capaldi group at
UCSF for a possible joint validation study positioning Tide as
the at-home companion to in-vault iSGRT; and (iv) a
cross-session anchor for inter-day plateau drift readout.

\section{Conclusions}

A browser-based smartphone application leveraging
\texttt{DeviceOrientationEvent} fused attitude can capture
sufficient signal to coach DIBH reproducibility at home, with
adaptive per-session calibration and position-match audio
feedback. Phone placement as a lever across the chest-belly seam
is essential to translate chest expansion into a usable rotational
signal. Tide complements rather than replaces in-clinic SGRT and
fills the practice gap between simulation and treatment.

\section*{Disclosures}

The authors report no conflicts of interest.

\section*{Sources of support}

Investigator-initiated. No external funding.

\section*{Declaration of generative AI assistance}

During the development of this application and preparation of this
manuscript, the authors used Anthropic's Claude (model: Opus 4.7,
1M context) as a coding assistant and writing aid. All algorithm
design decisions, sensor analyses, and clinical positioning were
reviewed and approved by the human authors.

\bibliographystyle{elsarticle-num}
\bibliography{refs}

\end{document}
